\newcommand*{\path5}{chapters/chapter5}

\section{Molecular Property Descriptors}

Molecular and physicochemical descriptors are very important in drug discovery.

%%%%%%%%%%%%%%%%%%%%%
\subsection{Molecular Weight (MW or MWt)}

The molecular weight (MW) is simply the sum of the masses of every heavy atom constituting the molecule (hydrogen is excluded). It is therefore easy to compute and very precise.

Large molecules have higher chances of containing undesirable moieties, therefore many drugs usually have a small molecular weight.

MW is used to compute the \emph{binding efficiency index}, which is the potency normalized by the MW. Added atoms should have a positive effect on the potency and not just increase the MW.


%%%%%%%%%%%%%%%%%%%%%%%%%%%%%%%%%%%%%%%%%%%%%%%%%%%%%%%%%%%%
\subsection{Octanol/Water Partition Coefficient (C$\log P$)}

The $\log P$ descriptor is widely used and is defined as the logarithm of the concentrations of a compound in two liquid phases (typically octanol and water):
\[
    \log P = \log \left(
        \frac{[\text{solute}]_{\ce{C8H18O}}}{[\text{solute}]_{\ce{H2O}}}
    \right)
\]

The $\log P$ coefficient indicates if a molecule is hydrophilic or hydrophobic, which affect the ADME (administration, absorption, transport and excretion) of such compound.

When $\log P$ is calculated instead of measured, it is usually called C$\log P$, but the method with which it has been computed should always be specified.


%%%%%%%%%%%%%%%%%%%%%%%%%%%%%%%%%%%%%%%%%%%%%%%%%%
\subsection{Topological Polar Surface Area (TPSA)}

The polar surface area (PSA) is the sum of the surface area of all polar atoms in a molecule (including connected hydrogens) and indicates the cell permeability of a molecule or its permeability of the blood–brain barrier (BBB).

The PSA can be computed from the 3D structure of the molecule, but this is computationally expensive. Therefore a topological method to compute the PSA has been developed, called topological PSA (TPSA).


%%%%%%%%%%%%%%%%%%%%%%%%%%%%%%%%%%%%%%%%%%%%%%%%%%%%%%%%%%%%%
\subsection{Hydrogen Bond Acceptors and Donors (HBA and HBD)}

A H-bond donor (HBD) is an lone pair acceptor, while a H-bond acceptor (HBA) is an lone-pair donor. HDA and HBD depend on the context of the atoms in the molecule and the pH ($pK_a$).

For simplicity, HBD is the sum of \ce{N-H} and \ce{OH} bonds while HBA is the sum of all nitrogen and oxygen atom.


%%%%%%%%%%%%%%%%%%%%%%%%%%%%%%%%%%%%
\subsection{Lipinski’s Rule-of-Five}

The majority of oral drugs are small and moderately lipophilic. Lipinski rule-of-five is an heuristic rule for \emph{druglikeness} (or oral \emph{bioavailability}):
\begin{itemize}
    \item Moleculear mass less than 500 Daltons
    \item $log P$ smaller than 5
    \item Maximum 5 HBD
    \item Maximum 10 HBA
\end{itemize}

The number of rotatable bonds (the sum of acyclic single bonds) is preferred to be smaller than 10.

